\documentclass[11pt,a4paper]{article}
\usepackage[a4paper,margin=2cm]{geometry}
\usepackage[utf8]{inputenc}
\usepackage[T1]{fontenc}
\usepackage[spanish,es-nodecimaldot,es-noquoting]{babel}
\usepackage{amsmath,amssymb,amsthm,booktabs,array,graphicx,multicol,pdflscape}
\usepackage{hyperref}
\hypersetup{colorlinks=true,linkcolor=blue!70!black,urlcolor=blue!70!black}
\setlength{\parskip}{4pt}
\title{\textbf{Taller de Cálculo Estocástico}\\large Sección 3: Integrales de Itô}
\author{Universidad Externado de Colombia}
\date{2026}
\begin{document}
\maketitle

\section*{Antes de empezar: herramientas que usaremos}
El browniano estándar \(B(t)\) empieza en cero. En un intervalo de duración \(t\), su cambio tiene distribución normal \(N(0,t)\), media cero y varianza \(t\). Por eso, si \(a\) es constante,
\[
\int_0^t a\,dB(u)=aB(t)\sim N(0,a^2t).
\]
Más generalmente, para una función determinista \(\sigma(u)\), la integral \(\int_0^t\sigma(u)\,dB(u)\) es normal, con media cero y varianza \(\int_0^t\sigma^2(u)\,du\). Una integral ordinaria \(\int_0^t\mu(u)\,du\) solo acumula deriva y no agrega azar. Esta distinción resuelve los ejercicios 1 y 2.

\textbf{Regla de Itô.} Si \(X_t\) satisface \(dX_t=a_t\,dt+b_t\,dB_t\), para una función \(f(t,x)\) suficientemente suave:
\[
df(t,X_t)=f_t\,dt+f_x\,dX_t+\tfrac12 f_{xx}(dX_t)^2,
\qquad (dB_t)^2=dt,\quad dt\,dB_t=(dt)^2=0.
\]
El término de segunda derivada es la corrección que distingue Itô del cálculo ordinario. En los ejercicios siguientes veremos cómo aparece.

\section*{Punto 1. Liquidez con deriva y volatilidad constantes}
\textbf{Enunciado.} La liquidez, en millones de dólares, es
\[
X(t)=X(0)+\int_0^t0{,}4\,dB(u)+\int_0^t0{,}1\,du,\qquad X(0)=2.
\]
Halle la distribución a 3, 6 y 12 meses y la probabilidad de liquidez negativa.

\subsection*{Qué datos tenemos y qué buscamos}
La parte aleatoria tiene coeficiente constante \(0{,}4\); su varianza por mes es \(0{,}4^2=0{,}16\). La deriva agrega \(0{,}1\) millones por mes. Usamos
\[
\int_0^t0{,}4\,dB(u)=0{,}4B(t),\qquad
\int_0^t0{,}1\,du=0{,}1t.
\]
Luego \(X(t)=2+0{,}1t+0{,}4B(t)\). La suma de una constante y una normal sigue siendo normal:
\[
X(t)\sim N\bigl(2+0{,}1t,\;0{,}16t\bigr).
\]
Aquí escribimos \(N(\text{media},\text{varianza})\), no \(N(\text{media},\text{desviación})\).

\subsection*{Cálculos y probabilidad negativa}
Para cada horizonte calculamos media \(m_t=2+0{,}1t\), varianza \(v_t=0{,}16t\), desviación \(\sqrt{v_t}\), y luego estandarizamos cero:
\[
\mathbb{P}(X(t)<0)=\Phi\left(\frac{0-m_t}{\sqrt{v_t}}\right),
\]
donde \(\Phi\) es la función de distribución normal estándar.
Por ejemplo, a 3 meses \(X(3)\sim N(2{,}30,0{,}48)\), y \(z=(0-2{,}30)/\sqrt{0{,}48}=-3{,}320\), cuya cola izquierda es aproximadamente \(0{,}00045\).

\textbf{Conclusión.} La liquidez esperada sube linealmente con el tiempo y la incertidumbre también crece. Aunque la probabilidad de saldo negativo es baja, aumenta entre 3 y 12 meses porque la desviación se acumula.

\section*{Punto 2. Deriva y volatilidad por tramos}
\textbf{Enunciado.} \(X(0)=1{,}8\) millones y
\[
X(t)=X(0)+\int_0^t\sigma(u)\,dB(u)+\int_0^t\mu(u)\,du,
\]
donde \(\sigma=0{,}30,\mu=0{,}06\) hasta el mes 6 y \(\sigma=0{,}50,\mu=0{,}12\) entre los meses 6 y 12. Halle ley y probabilidad de saldo negativo en \(t=3,6,12\).

\subsection*{Paso 1: separar lo aleatorio de lo acumulado}
Para una deriva que cambia con el tiempo, la media es \(m(t)=X(0)+\int_0^t\mu(u)\,du\). Para la integral browniana, la varianza es \(v(t)=\int_0^t\sigma^2(u)\,du\). El cuadrado en \(\sigma^2\) es esencial.
\[
X(t)\sim N\left(1{,}8+\int_0^t\mu(u)\,du,\;\int_0^t\sigma^2(u)\,du\right).
\]
Si \(t\le6\), las integrales son \(0{,}06t\) y \(0{,}30^2t=0{,}09t\). Si \(t>6\), se acumulan primero seis meses del tramo 1 y luego \(t-6\) meses del tramo 2:
\[
m(t)=1{,}8+0{,}06(6)+0{,}12(t-6),\quad
v(t)=0{,}09(6)+0{,}50^2(t-6).
\]

\subsection*{Paso 2: sustituir los horizontes}
A 12 meses, por ejemplo, la media es \(1{,}8+0{,}06(6)+0{,}12(6)=2{,}88\). La varianza es \(0{,}30^2(6)+0{,}50^2(6)=2{,}04\), y \(X(12)\sim N(2{,}88,2{,}04)\). Para la probabilidad negativa usamos \(z=-2{,}88/\sqrt{2{,}04}=-2{,}016\).

\textbf{Conclusión.} Al entrar al segundo tramo sube tanto la deriva como la volatilidad. La liquidez media crece, pero el aumento de varianza eleva la probabilidad de un saldo negativo al mes 12.

\section*{Puntos 4--7. Probabilidades de superar un precio bajo GBM}
\textbf{Enunciados.} Para cada acción se da el modelo
\[
S(t)=S_0\exp(\mu t+\sigma B(t)).
\]
Calcule la probabilidad de superar el nivel indicado.

\subsection*{Método común, explicado una vez}
Como \(B(t)\sim N(0,t)\), podemos escribir \(B(t)=\sqrt{t}\,Z\), con \(Z\sim N(0,1)\). Tomamos logaritmo para convertir el umbral de precio en un umbral normal:
\[
S(t)>K
\Longleftrightarrow
\ln S_0+\mu t+\sigma\sqrt{t}\,Z>\ln K
\Longleftrightarrow
Z>\frac{\ln(K/S_0)-\mu t}{\sigma\sqrt{t}}.
\]
Llamemos \(z_K\) al último cociente. La respuesta es \(1-\Phi(z_K)\). En cada inciso sustituimos primero los datos y luego buscamos la cola derecha en la normal estándar.

\subsection*{Punto 4}
Datos: \(S_0=50\), \(K=80\), \(\mu=0{,}12\), \(\sigma=0{,}30\), \(t=2\) años.
\[
z_K=\frac{\ln(80/50)-0{,}12(2)}{0{,}30\sqrt2}
=\frac{0{,}4700-0{,}2400}{0{,}4243}=0{,}5421.
\]
Por tanto, \(\mathbb{P}(S(2)>80)=1-\Phi(0{,}5421)\approx\boxed{0{,}2939}\), es decir, \(29{,}39\%\).

\subsection*{Punto 5}
Siete meses son \(t=7/12=0{,}5833\) años. Con \(S_0=25\), \(K=30\), \(\mu=0{,}16\), \(\sigma=0{,}20\):
\[
z_K=\frac{\ln(30/25)-0{,}16(7/12)}{0{,}20\sqrt{7/12}}
=\frac{0{,}1823-0{,}0933}{0{,}1528}=0{,}5826.
\]
Así, \(\mathbb{P}(S(7/12)>30)=1-\Phi(0{,}5826)\approx\boxed{0{,}2801}\) (28,01\%).

\subsection*{Punto 6}
Con \(S_0=60\), \(K=90\), \(\mu=0{,}08\), \(\sigma=0{,}25\) y \(t=1{,}5\):
\[
z_K=\frac{\ln(90/60)-0{,}08(1{,}5)}{0{,}25\sqrt{1{,}5}}
=\frac{0{,}4055-0{,}1200}{0{,}3062}=0{,}9323.
\]
Entonces \(\mathbb{P}(S(1{,}5)>90)=1-\Phi(0{,}9323)\approx\boxed{0{,}1756}\) (17,56\%).

\subsection*{Punto 7}
Con \(S_0=45\), \(K=70\), \(\mu=0{,}10\), \(\sigma=0{,}40\), \(t=1\):
\[
z_K=\frac{\ln(70/45)-0{,}10}{0{,}40}
=\frac{0{,}4418-0{,}1000}{0{,}40}=0{,}8546.
\]
Luego \(\mathbb{P}(S(1)>70)=1-\Phi(0{,}8546)\approx\boxed{0{,}1964}\) (19,64\%).

\subsection*{Conclusión de los puntos 4--7}
En los cuatro ejercicios el paso clave es pasar de \(S(t)>K\) a una desigualdad para una normal estándar. Convertir meses a años antes de sustituir evita mezclar unidades de tiempo con parámetros anualizados.

\section*{Punto 8. Aplicaciones de Itô}
Para una función (f(t,B(t))), usamos (df=f_tdt+f_xdB+\tfrac12f_{xx}dt). En (tB(t)), las derivadas son (f_t=B, f_x=t, f_{xx}=0), así que (d(tB)=Bdt+t,dB). Para (B^2), (f'=2B,f''=2), y (d(B^2)=2B,dB+dt). Para (B^3), (f'=3B^2,f''=6B), y (d(B^3)=3B^2dB+3Bdt).

\section*{Punto 9. Regla de producto}
Itô usa (d(XY)=X,dY+Y,dX+dX,dY). Para un factor determinista, el término cruzado con (dB) es cero porque (dt,dB=0).
\[
d(t^2B)=t^2dB+2tBdt,\qquad d(e^{\alpha t}B)=e^{\alpha t}dB+\alpha e^{\alpha t}Bdt.
\]
Si (J(t)=\int_0^tB(s)ds), entonces (dJ=Bdt), y
\[
d(BJ)=JdB+B,dJ+dB,dJ=JdB+B^2dt.
\]

\section*{Punto 10. Fórmula de Itô}
Para (f(B)), (df=f'(B)dB+\tfrac12f''(B)dt). Aplicamos derivadas ordinarias:
\[
f(x)=x^4:\quad f'=4x^3,\ f''=12x^2\quad\Longrightarrow\quad d(B^4)=4B^3dB+6B^2dt.
\]
\[
f(x)=\sin x:\quad f'=\cos x,\ f''=-\sin x\quad\Longrightarrow\quad d(\sin B)=\cos BdB-\tfrac12\sin Bdt.
\]
\[
f(x)=e^{\beta x}:\quad f'=\beta e^{\beta x},\ f''=\beta^2e^{\beta x}\quad\Longrightarrow\quad d(e^{\beta B})=\beta e^{\beta B}dB+\tfrac12\beta^2e^{\beta B}dt.
\]

\section*{Punto 11. Esperanza y varianza}
(Y=B(t)^2+2t). Puesto que (B(t)\sim N(0,t)), (E[B^2]=t) y (E[B^4]=3t^2):
\[
E[Y]=t+2t=3t,\qquad \operatorname{Var}(Y)=E[B^4]-E[B^2]^2=3t^2-t^2=2t^2.
\]
Comprobación con Itô: (dY=2B,dB+3dt). La integral aleatoria tiene esperanza cero; su varianza es (E[\int_0^t4B(s)^2ds]=4\int_0^t sds=2t^2).

\section*{Punto 12. Cambio del orden de integración}
Probamos (\int_0^tB(s)ds=\int_0^t(t-s)dB(s)). Aquí (t) es constante, un extremo fijo; (s) es la variable que recorre la integral. Por eso (t-s) es una función determinista de (s).
Como (B(s)=\int_0^s dB(u)),
\[
\int_0^tB(s)ds=\int_0^t\left(\int_0^s dB(u)\right)ds.
\]
La región es (0\le u\le s\le t). Invirtiendo el orden, (u) recorre de 0 a (t), y para cada (u), (s) va de (u) a (t):
\[
\int_0^t\int_0^s dB(u)ds=\int_0^t\left(\int_u^t ds\right)dB(u)=\int_0^t(t-u)dB(u).
\]
Renombrando la variable muda (u) como (s), obtenemos la igualdad. La pizarra la expresa también (tB(t)-\int_0^t s,dB(s)), forma equivalente por integración por partes.

\section*{Punto 13. Integral de (B^2dB)}
La extracción del enunciado parece perder el cuadrado: la identidad de la pizarra es con (B(s)^2dB(s)). Sin él, la igualdad indicada no sería correcta. Partimos de Itô para el cubo:
\[
d(B^3)=3B^2dB+3Bdt.
\]
Integramos desde 0 hasta (t), usamos (B(0)=0), y despejamos:
\[
B(t)^3=3\int_0^tB(s)^2dB(s)+3\int_0^tB(s)ds,
\quad\boxed{\int_0^tB(s)^2dB(s)=\tfrac13B(t)^3-\int_0^tB(s)ds}.
\]

\section*{Punto 14. Taylor para funciones de (t,W(t))}
Taylor hasta segundo orden da (df=f_tdt+f_WdW+\tfrac12f_{WW}dt): el cuadrado ((dW)^2) se reemplaza por (dt), y los términos (dt,dW), (dt^2) desaparecen. Interpretamos las expresiones como hallar sus diferenciales.
Para (f_1=2tW+W^2):
\[
(f_1)_t=2W,\quad(f_1)_W=2t+2W,\quad(f_1)_{WW}=2,
\]
\[
df_1=2Wdt+(2t+2W)dW+\tfrac12(2)dt
=\boxed{(2t+2W)dW+(2W+1)dt}.
\]
Para (f_2=W^2+3tW^3):
\[
(f_2)_t=3W^3,\quad(f_2)_W=2W+9tW^2,\quad(f_2)_{WW}=2+18tW.
\]
\[
\begin{aligned}
df_2&=3W^3dt+(2W+9tW^2)dW+\tfrac12(2+18tW)dt\
&=\boxed{(2W+9tW^2)dW+(1+9tW+3W^3)dt}.
\end{aligned}
\]
La segunda derivada respecto a (W) produce la corrección adicional con (dt).

\section*{Punto 15. Ecuación diferencial del GBM}
Partimos de (S_t=S_0e^{\mu t+\sigma B(t)}) y definimos (f(t,x)=S_0e^{\mu t+\sigma x}). Sus derivadas son (f_t=\mu f, f_x=\sigma f, f_{xx}=\sigma^2f). Por Itô:
\[
dS_t=\mu S_tdt+\sigma S_tdB+\tfrac12\sigma^2S_tdt,
\quad\boxed{dS_t=(\mu+\tfrac12\sigma^2)S_tdt+\sigma S_tdB(t)}.
\]
El término adicional (\tfrac12\sigma^2) aparece por la segunda derivada de la exponencial.

\section*{Punto 16. Producto de factor determinista y proceso OU}
(X_t=e^{kt}), (dX_t=ke^{kt}dt); (Y_t) satisface (dY_t=k(\Theta-Y_t)dt+\sigma dB(t)), (Y_0=r_0). Sea (Z_t=X_tY_t=e^{kt}Y_t).
\subsection*{Diferencial del producto}
\[
dZ_t=X_t dY_t+Y_t dX_t+dX_t dY_t.
\]
El último término es cero porque (dX_t) contiene solo (dt), y (dt^2=dt,dB=0). Sustituyendo:
\[
\begin{aligned}
dZ_t&=e^{kt}[k(\Theta-Y_t)dt+\sigma dB]+kY_te^{kt}dt\
&=k\Theta e^{kt}dt+\sigma e^{kt}dB(t).
\end{aligned}
\]
Los términos con (Y_t) se cancelan.
\subsection*{Integración y distribución}
Con (Z_0=r_0), integramos:
\[
Z_t=r_0+k\Theta\int_0^t e^{ks}ds+\sigma\int_0^t e^{ks}dB(s).
\]
Para (k\ne0), la primera integral es ((e^{kt}-1)/k). La integral estocástica tiene media cero y varianza \(\sigma^2\int_0^t e^{2ks}ds\):
\[
\boxed{e^{kt}Y_t=r_0+\Theta(e^{kt}-1)+\sigma\int_0^t e^{ks}dB(s)},
\]
\[
\boxed{E[e^{kt}Y_t]=r_0+\Theta(e^{kt}-1)},\qquad
\boxed{\operatorname{Var}(e^{kt}Y_t)=\sigma^2\frac{e^{2kt}-1}{2k}}.
\]
La distribución es normal con esa media y varianza. Si (k=0), los límites son media (r_0) y varianza \(\sigma^2t\).

\newpage
\section*{Formulario de la sección 3}
\scriptsize
\setlength{\columnsep}{0.8cm}
\begin{multicols}{2}
\textbf{Integrales y normalidad}
\[
\int_0^t c\,dB=cB(t),\quad E[\int_0^t h\,dB]=0,\quad
\operatorname{Var}(\int_0^t h\,dB)=\int_0^t h^2ds.
\]
\[
X(t)\sim N(X_0+\int_0^t\mu ds,\int_0^t\sigma^2ds),\quad
P(X<0)=\Phi(-E[X]/\sqrt{\operatorname{Var}(X)}).
\]
\textbf{GBM}
\[
S_t=S_0e^{\mu t+\sigma B_t},\quad
P(S_t>K)=1-\Phi((\ln(K/S_0)-\mu t)/(\sigma\sqrt t)).
\]
\[dS_t=(\mu+\tfrac12\sigma^2)S_tdt+\sigma S_tdB_t.\]
\textbf{Reglas de Itô}
\[df=f_tdt+f_xdB+\tfrac12f_{xx}dt,\quad(dB)^2=dt,\quad dt,dB=(dt)^2=0.\]
\[d(XY)=X,dY+Y,dX+dX,dY.\]
\[d(B^2)=2B,dB+dt,\quad d(B^3)=3B^2dB+3Bdt.\]
\[d(B^4)=4B^3dB+6B^2dt.\]
\[d(\sin B)=\cos B,dB-\tfrac12\sin Bdt.\]
\[d(e^{\beta B})=\beta e^{\beta B}dB+\tfrac12\beta^2e^{\beta B}dt.\]
\textbf{Integrales útiles}
\[\int_0^tB(s)ds=\int_0^t(t-s)dB(s).\]
\[\int_0^tB(s)^2dB(s)=\tfrac13B(t)^3-\int_0^tB(s)ds.\]
\textbf{Momentos}
\[E[B_t]=0,\quad E[B_t^2]=t,\quad E[B_t^4]=3t^2.\]
\[E[B_t^2+2t]=3t,\quad\operatorname{Var}(B_t^2+2t)=2t^2.\]
\textbf{Producto OU}
\[d(e^{kt}Y_t)=k\Theta e^{kt}dt+\sigma e^{kt}dB_t.\]
\[E[e^{kt}Y_t]=r_0+\Theta(e^{kt}-1),\quad\operatorname{Var}(e^{kt}Y_t)=\sigma^2(e^{2kt}-1)/(2k).\]
\end{multicols}
\normalsize
\end{document}
